\section{观察时域与历史信息的影响}
\label{sec:records}
定理\ref{thm:main}通过下一时刻的事件反映当前偏差，并利用协方差标签校准正常静默概率。
为分析这两类信息的作用，分别考察截至目标时刻的完整日志和不含历史条件的下一事件。
\begin{proposition}[当拍截止记录的局限]
\label{prop:deletion}
在主定理的任意固定 $n$ 维模型下，若只约束截至目标时刻的完整日志 $H_k$，
则充分小的 $b>0$ 时仍存在满足峰值约束的因果攻击，使 $H_k$ 与无攻击运行同律，
同时产生 $\Theta(b^2)$ 的正损伤。
\end{proposition}
\begin{proof}
目标时刻前不攻击。目标时刻发送后，对真实创新的发送条件密度施加步长为 $\pm bv$ 的一步
Metropolis--Hastings核\cite{metropolis1953,hastings1970}。
详细平衡保证收到的载荷保持原分布，前史及当前发送事件也不变，故 $H_k$ 同律。
接受的非零移动引入估计偏差，其损伤具有正的二次下界；完整推导见附录\ref{app:mh}。
\end{proof}

命题\ref{prop:deletion}表明，当拍日志保持正常分布，并不能保证接收估计与真实状态之间的误差不变。
该偏差经预测反馈影响下一次触发，因此包含下一事件的记录仍受定理\ref{thm:main}约束。
若保留下一事件，但删去全部历史条件，则需要进一步检验不同历史的概率变化能否抵消。记
$\Delta_{\rm last}=\TV(\Law_\pi(\gamma_{k+1}),\Law_0(\gamma_{k+1}))$。
\begin{proposition}[末事件位的标量反例]
\label{prop:last-bit}
取标量模型 $A=1$、$P_0=4/5$、$Q=2$、$R=1/5$、$Z=2$ 及目标时刻 $k=1$。
充分小的 $b>0$ 下存在满足峰值约束的因果策略，使
\begin{equation}
 \begin{gathered}
 \Delta_{\rm last}=0,\quad J_1=c_*b^2+O(b^4),\\
 c_*\approx0.00049907284045>0.
 \end{gathered}
 \label{eq:last-counterexample}
\end{equation}
\end{proposition}
\begin{proof}
构造使末次静默概率分别增加和减小的两种合法策略，再以依赖给定幅值 $b$ 的固定私有概率混合，精确抵消该概率差。
混合后的二次损伤系数严格为正；策略、混合权重及系数的有理区间验证见附录\ref{app:last-bit}。
\end{proof}
在命题\ref{prop:last-bit}的标量模型中，不同历史下的概率变化抵消后，末事件位仍与正常运行同律，但真实损伤为正。
因此，定理\ref{thm:main}不能在全部模型上推广至仅含下一事件边缘分布的约束。

精确协方差标签 $P_k^+$ 为下一事件提供了相应的正常概率基准。
它由事件历史确定；若该映射为单射，两者具有相同的统计区分信息。
精确标签由此具有明确的概率解释，但不必然减少历史记录的信息量；量化后能否保留损伤保证仍需另行分析。
此外，本文的TV约束针对总体分布。将其用于统计认证时，还需计入概率估计误差。
实值得分、随机二值输出及独立重置实验下的有限样本置信界见附录\ref{app:statistical}。
